\section{适定性与能量方法}
\label{sec:15-Differential-Equations-and-Dynamical-Systems/03-PDE-Introduction/03}

手里有一张失焦的照片，想把它复原。失焦这个操作在数学上就是让图像沿时间跑一段热方程——高频细节被抹掉，边缘变糊。既然正向过程有明确的方程，反过来解一段不就复原了？把模糊图往回推的代码写起来只要几行，跑出来的结果却是一屏幕高频噪声，原图的影子都看不见。代码没写错，方程也没写错。问题出在这个反问题本身。

ODE 那边，存在唯一性由一句 Lipschitz 条件管住（见\hyperref[sec:15-Differential-Equations-and-Dynamical-Systems/01-ODE-Theory/03]{``存在唯一性为什么需要 Lipschitz''}）。PDE 没有这样一网打尽的定理，取而代之的是 Hadamard 提出的三条标准：解存在、解唯一、解连续依赖于数据。前两条读起来天经地义，第三条最容易被当成技术性附注跳过，而失焦复原失败的原因恰恰全在第三条。

\subsection{第三条不成立时，误差不是变大，是变成主角}

把热方程 \(\partial_tu=\partial_{xx}u\) 的时间方向掉个头，等价于问：能不能从此刻的温度反推一小时前的温度。取一段高频扰动作为当前数据，\(u_0(x)=\frac1n\sin(nx)\)，从当前时刻 \(t=0\) 出发的解是

\[
u(t,x)=\frac1n\,e^{-n^2t}\sin(nx),
\qquad
u(-\tau,x)=\frac1n\,e^{n^2\tau}\sin(nx)\quad(\tau>0).
\]

往前推（\(t>0\)）时因子 \(e^{-n^2t}\) 让它迅速归零；往过去倒推 \(\tau\) 个时间单位则换成 \(e^{n^2\tau}\) 的放大。数据上一个幅度只有 \(1/n\) 的抖动，倒推同样一段时间之后幅度变成 \(\frac1n e^{n^2\tau}\)。令 \(n\to\infty\)：数据端的误差趋于零，答案端的误差趋于无穷。这两者之间没有任何一致的界，连续依赖彻底破产。

\begin{center}
\begin{tikzpicture}[>=stealth]
  \begin{scope}
    \draw[->] (0,0) -- (3.4,0) node[right] {\small \(x\)};
    \draw[thick] plot[domain=0:3.2,samples=120,variable=\x]
      ({\x},{1.1+0.5*sin(56.25*\x)});
    \draw[dashed] plot[domain=0:3.2,samples=200,variable=\x]
      ({\x},{1.1+0.5*sin(56.25*\x)+0.06*sin(720*\x)});
    \node at (1.6,2.55) {\small 观测数据：两条几乎重合};
  \end{scope}

  \begin{scope}[shift={(5.0,0)}]
    \draw[->] (0,0) -- (3.4,0) node[right] {\small \(x\)};
    \draw[thick] plot[domain=0:3.2,samples=120,variable=\x]
      ({\x},{1.1+0.5*sin(56.25*\x)});
    \draw plot[domain=0:3.2,samples=500,variable=\x]
      ({\x},{1.1+0.5*sin(56.25*\x)+0.68*sin(720*\x)});
    \node at (1.6,2.55) {\small 倒推结果：面目全非};
  \end{scope}
\end{tikzpicture}
\end{center}

\emph{肉眼分不开的两组数据，倒着解一段热方程后分道扬镳；放大倍数只由噪声的频率决定。}

图里那点几乎看不见的虚线抖动，正是任何真实观测都躲不掉的传感器噪声或量化误差。它在正向问题里无关痛痒，在反向问题里被 \(e^{n^2\abs{t}}\) 挑出来放到最前面。上一节说抛物方程不能反演，定量版本就写在这个指数里：正向的抹平机制被逆转成高频爆炸，而噪声天然富含高频。

同样的把戏对 Laplace 方程的``初值问题''也有效。若对 \(\partial_{xx}u+\partial_{yy}u=0\) 沿 \(y\) 方向给出初始数据，高频分量按 \(e^{ny}\) 增长，结论一样。Hadamard 当年拿这类例子说明的正是：三条标准缺一条的问题，在物理上和计算上都无法交付。

\subsection{放大倍数就是条件数，只是它现在没有上界}

这件事在数值分析里已经有名字。\hyperref[sec:08-Numerical-Analysis/01-Floating-Point-and-Error/03]{``条件数衡量问题本身的敏感性''}把相对误差的放大倍数定义为条件数，并强调它属于问题而非算法——条件数大的问题，再稳定的算法也救不回精度。反解热方程把这套语言推到极端：离散之后的算子在频率 \(k\) 上的奇异值是 \(e^{-\kappa k^2t}\)，求逆就是除以它，条件数随网格加密指数增长，没有上界可言。换更好的求解器、用双精度、加迭代次数，都不解决问题。

实践中的出路是承认这一点，然后改问题。给目标函数加一项 \(\lambda\norm{\nabla u}^2\) 或 \(\lambda\norm{u}^2\)，即 Tikhonov 正则化，等于人为给小奇异值设一个下限，把无界的放大截断在 \(1/\lambda\) 的量级；代价是分辨率——被截掉的那部分高频，本来就已经被正向过程消灭得低于噪声水平，任何方法都无法诚实地恢复。去模糊、CT 重建、地震反演，做法都是这一套。扩散模型的逆向采样也不去倒解前向方程，而是另训一个网络来估计得分（见\hyperref[sec:06-Random-Processes/06-Diffusion-and-SDE/05]{``得分匹配：怎样从样本学习得分''}），用学到的先验补上被抹掉的信息。

\subsection{能量方法用一条不等式同时买下唯一与稳定}

正向问题这边，好消息是不需要知道解的公式也能把适定性论证清楚。能量方法是 PDE 里最接近 Lyapunov 函数的工具（见\hyperref[sec:15-Differential-Equations-and-Dynamical-Systems/02-Stability-and-Dynamical-Systems/02]{``稳定性与 Lyapunov 函数''}）：造一个非负泛函，证明它沿解不增。

以齐次 Dirichlet 边界下的热方程为例，取

\[
E(t)=\frac12\int_0^L u(t,x)^2\,dx .
\]

对 \(t\) 求导，把方程代进去，再对空间分部积分。边界项因为 \(u\) 在两端为零而消失：

\[
\begin{aligned}
E'(t)&=\int_0^L u\,\partial_tu\,dx
=\kappa\int_0^L u\,\partial_{xx}u\,dx\\
&=\kappa\Bigl[u\,\partial_xu\Bigr]_0^L-\kappa\int_0^L(\partial_xu)^2\,dx
=-\kappa\int_0^L(\partial_xu)^2\,dx\le0 .
\end{aligned}
\]

骨架就这么长，值得对账的是用掉了哪些条件。求导与积分交换需要解有足够正则性，分部积分需要 \(\partial_xu\) 在 \([0,L]\) 上平方可积，边界项归零需要齐次 Dirichlet 条件——换成齐次 Neumann 条件（\(\partial_xu=0\)）边界项照样消失，结论不变；换成非齐次边界，边界项就是外界注入的能量，不等式要相应改写。\(\kappa>0\) 也是必需的，符号一反整个论证倒过来，这正是上一小节那个反问题的来处。

唯一性由此立刻到手。设 \(u\) 与 \(v\) 是同一组数据下的两个解，令 \(w=u-v\)，则 \(w\) 满足齐次方程、零初值、零边界，于是 \(E_w(0)=0\) 且 \(E_w\) 不增，非负泛函被夹死为零，\(w\equiv0\)。连续依赖同理：两组初值的 \(L^2\) 距离控制住此后所有时刻解的 \(L^2\) 距离，误差不会自己长大。再补一条 Poincaré 不等式 \(\int u^2\le C\int(\partial_xu)^2\)（齐次边界下成立），可得 \(E'\le-2\kappa C^{-1}E\)，能量指数衰减——这是``抹平''的定量说法。

波动方程有结构相同的能量，只是配方里多了动能项：

\[
E(t)=\frac12\int_0^L\Bigl[(\partial_tu)^2+c^2(\partial_xu)^2\Bigr]dx .
\]

同样求导、代方程、分部积分，两项恰好相消，\(E'(t)=0\)。能量守恒既给出唯一性，也解释了波形为什么不衰减。数值格式好不好用，一条快速检验就是看它离散之后这个量是否也守恒——耗散型格式会让波幅慢慢矮下去，那是格式的黏性，不是物理。

\subsection{存在性得靠构造，估计本身不生产解}

能量估计的前提是``设解存在''，它管的是唯一与稳定，产不出解。存在性要另找路子：分离变量把解写成级数（下一节），半群理论把演化写成算子指数，或者先离散再让网格趋于零——一个收敛的数值格式本身就是一份构造性的存在证明草稿。

对本节涉及的三个线性原型，在通常的光滑性与相容性要求下，初边值问题都是适定的。适定性出问题的场合反倒有规律可循：时间方向弄反，数据类型与方程类型不匹配，边界条件给多了或给少了。三种情形分别对应第三条、第一条、前两条的失败。

至此，``解''这个字仍默认是逐点满足方程的经典解。下一节先把这类解在规则区域上完整地算出来，第五节再回头追问：当解本身就不可微时，方程该怎么写。
